%============================================================
\chapter{The Supplier of the Information Structure: Settlement by Protocol}
\label{ch:protocol}
%============================================================

\section{The premises imposed in this chapter}

Chapter~\ref{ch:solo} imposed constraints on resources and Chapter~\ref{part:unmanned} took
their limit. Both take the information structure $\mathcal{G}$ of Chapter~\ref{sec:info} as
given. This chapter varies $\mathcal{G}$ itself and asks
\textbf{how the set of feasible $\Phi$ changes when the supplier of $\mathcal{G}$ is
replaced}.

Two things should be said first.

First, \textbf{this chapter is deduction, not evidence.} The validity of the conclusions
depends on the validity of the premises.

Second, \textbf{nothing here is new.} Every relation stated has a counterpart in the existing
literature, and a correspondence table is given in Appendix~\ref{app:priorwork}. The purpose
of the chapter is to show how the theory built in Part~\ref{part:theory} works under
specific conditions.

What is replaced is the $\mathcal{G}$ of Chapter~\ref{sec:info}; the question is who
guarantees the measurability condition of Proposition~\ref{prop:implementable}.

\begin{definition}[Protocol]
\label{def:protocol}
A set of rules published in advance, whose execution every participant can verify
independently, is called a \emph{protocol} and written $\mathcal{P}$. The tuple of variables
$\mathcal{P}$ holds is called its \emph{state}. State transitions are taken to be atomic:
no part of a transition is applied on its own.
\end{definition}

The premises are as follows.

\begin{enumerate}[label=(\arabic*)]
\item Settlement $\pi$ is expressed as a state transition on $\mathcal{P}$
\item The execution of $\mathcal{P}$ can be verified independently by every participant (Definition~\ref{def:protocol})
\item Whether delivery $\delta$ is contained in the state of $\mathcal{P}$ depends on the object
\end{enumerate}

The implementation in wide use that satisfies premises (1) and (2) is settlement on a public
distributed ledger, that is, a blockchain. Not everything called by that name satisfies the
premises, however. Where verifiers are restricted, or where the state is held by a particular
operator, (2) fails and the deductions of this chapter do not apply.

Premises (1) and (2) are a relaxation. The $\mathcal{G}=\bigcap_i\mathcal{G}_i$ of
Chapter~\ref{sec:info} is normally supplied by a third party outside the transacting
parties, on whom verification and enforcement of $\pi$ depended. Premises (1) and (2) make
that unnecessary.

\textbf{Only (3) cannot be fixed as a premise}, and that is what forces this chapter's case
split.

\begin{center}
\begin{tabularx}{\textwidth}{@{}llX@{}}
\toprule
\multicolumn{3}{@{}l}{\textbf{Theoretical quantities in this chapter}}\\
\midrule
$\mathcal{G}$ & Chapter~\ref{sec:info} & the object whose supplier is replaced \\
The index $i$ of $\kappa_i$ & Eq.~\eqref{eq:kappa} & becomes an identifier on the protocol. \textbf{additivity is lost} \\
$\dbar-\delta$ & Definition~\ref{def:breakage-def} & observable as protocol state in the ``inside'' partition \\
$E$ & Eq.~\eqref{eq:cash} & appears as collateral \\
Degrees of freedom (2), (3) & Section~\ref{sec:dof} & the sieve cuts along these two \\
\bottomrule
\end{tabularx}
\captionof{table}{Theoretical quantities in this chapter: the symbols involved when the supplier of the information structure is replaced.}
\label{tab:protocol-quantities}
\end{center}

\section{Translating the premises}

\subsection{Three functions the third party performed}

The role of the supplier of $\mathcal{G}$ in a contract divides into three.

\begin{center}
\begin{tabularx}{\textwidth}{@{}lXl@{}}
\toprule
Function & Content & Usual bearer \\
\midrule
(a) Recording & holding the state of $\kappa_i$ & banks, payment processors \\
(b) Verification & judging whether the condition triggering $\pi$ has occurred & courts, auditors \\
(c) Enforcement & transferring value when the condition holds & courts, banks \\
\bottomrule
\end{tabularx}
\captionof{table}{Three functions a third party performs in a contract.}
\label{tab:third-party-functions}
\end{center}

By premises (1) and (2), $\mathcal{P}$ can always bear (a). It can bear (b) so long as the
condition can be judged from the state of $\mathcal{P}$, and (c) so long as what is
transferred is the state of $\mathcal{P}$.

\subsection{A case split by where delivery resides}

Premise (3) produces two cases.

\begin{center}
\begin{tabularx}{\textwidth}{@{}llXX@{}}
\toprule
Partition & Where $\delta$ resides & What enters $\mathcal{G}$ & Functions $\mathcal{P}$ can bear \\
\midrule
Inside & the state of $\mathcal{P}$ & both $\delta$ and $\pi$ & (a), (b), (c) \\
Outside & outside $\mathcal{P}$ & only $\pi$ & (a), and (b), (c) concerning $\pi$ \\
\bottomrule
\end{tabularx}
\captionof{table}{The partition by where delivery resides, and the functions the protocol can bear.}
\label{tab:protocol-cases}
\end{center}

The difference between the partitions bears directly on degree of freedom (1) of
Section~\ref{sec:dof}, timing, and fixes the sign of $\kappa_i$ in~\eqref{eq:kappa}.

\begin{proposition}[Simultaneity]
\label{prop:simultaneity}
In the ``inside'' partition, for any $\Phi$ there is an implementation with $\kappa_i=0$.
In the ``outside'' partition, $\kappa_i=0$ cannot be achieved for a $\Phi$ that involves
$\delta$.
\end{proposition}

\begin{proof}
In the ``inside'' partition both $\delta$ and $\pi$ are state transitions of $\mathcal{P}$.
By the atomicity of Definition~\ref{def:protocol} the two can be composed as a single
transition,
and $D_i$ and $P_i$ then move together, so the $\kappa_i$ of~\eqref{eq:kappa} is identically
$0$.

In the ``outside'' partition $\delta$ is not a state of $\mathcal{P}$, so $\mathcal{P}$ has
no time for $\delta$. $\kappa_i=0$ requires $\tau_i=0$, but $\tau_i$ is not determined on
$\mathcal{G}$.
\end{proof}

That is, \textbf{in the ``inside'' partition degree of freedom (1) becomes a full design
variable, while in the ``outside'' partition $\kappa\neq 0$ is forced}.
Immediate exchange in the real world achieves simultaneity because the parties observe each
other at the same moment; $\mathcal{P}$ has no such observation. This is why 1-1 (immediate
exchange) disappears in the ``outside'' partition in Section~\ref{sec:protocol-filter}.

\begin{remark}[The partition is a property of $\Phi$, not of the object]
\label{rem:case-is-phi}
For one and the same object the partition changes with what is taken to be delivery. A $\Phi$
exchanging a protocol balance for another balance is inside; a $\Phi$ exchanging the same
balance for currency outside the protocol is outside.

That a record of ownership sits on the protocol does not make the partition inside. Where a
non-fungible token is tied to a physical artwork or a membership, what moves is the record and
the delivery is a handover in the real world. Because $\delta$ lies outside $\mathcal{P}$,
the partition is outside. That the two cannot be identified shows up as the fact that a
change of record-holder does not cause a handover. For the same token, a $\Phi$ whose
delivery is only the exercise of a right on the protocol is inside.

The question ``is that object inside or outside?'' therefore does not arise. What can be
asked is ``is that $\Phi$ inside or outside?'' This has the same form as the conclusion of
Section~\ref{sec:unit}: here too the unit of observation must be $\Phi$ rather than the
object.
\end{remark}

\begin{remark}[Intermediation in the ``outside'' partition]
\label{rem:oracle}
To bring a condition concerning $\delta$ into $\mathcal{P}$ in the ``outside'' partition, some
party must write an external fact into the state. What $\mathcal{P}$ can verify is only that
a prescribed signature is attached; it does not verify that the value written is factual.
That party therefore fails premise (2), and \textbf{the third party supposedly removed
returns}.

The writing party may be a single entity or a quorum of several, so removal is continuous
rather than discrete. This chapter does not treat that continuum and handles only the two
ends (Section~\ref{sec:protocol-limits}). A classification of implementations is in
\cite{pasdar2023}.
\end{remark}

\subsection{The binding constraint switches}

In Chapter~\ref{ch:solo}, which of the financial and capacity constraints binds first
switches with the premises. The same happens here.

Because premises (1) and (2) widen $\mathcal{G}$, the measurability condition of
Proposition~\ref{prop:implementable} is looser in the ``inside'' partition than anywhere else
in this text. \textbf{$\mathcal{G}$ therefore ceases to bind and something else does.}
Section~\ref{sec:collateral} identifies what.

\section{Identifiers diverge from parties}
\label{sec:identifier}

This section is a by-product, placed where Chapter~\ref{ch:solo} recorded the correspondence
$A\to\kappa_S<0$ for the cloud.

The $\kappa_i$ of~\eqref{eq:kappa} is defined per party $i$ and summed as $\sum_i\kappa_i$
in~\eqref{eq:cash}. That summation presumes \textbf{that the index denotes a party uniquely}.

What can be identified on a protocol is an identifier. $\mathcal{P}$ cannot rule out one
party holding several identifiers, because generating an identifier costs nothing. This is a
known problem \cite{douceur2002}.

\begin{remark}[Loss of additivity]
\label{rem:additivity}
Writing $j$ for an identifier, $i$ for a party, and $j\mapsto\iota(j)$ for the
correspondence, what is observed is $\kappa_j$ while what is needed is
\begin{equation}
\kappa_i = \sum_{j:\,\iota(j)=i} \kappa_j
\label{eq:kappa-aggregation}
\end{equation}
Since $\iota$ is not observable, \eqref{eq:kappa-aggregation} cannot be evaluated.

\textbf{This is a limit of the definitions of Chapter~\ref{sec:phi}, not a consequence of
this chapter's premises.} It shows that a setting which fixes $N$ and ties the index to a
party does not fit an environment where identifiers are freely created. It originates in the
same place as the circulation of claims in Section~\ref{sec:token-limit}, namely the setting
of the domain. The theory is not amended.
\end{remark}

Remark~\ref{rem:additivity} bears directly on family 7.

\begin{proposition}[A $\Phi$ requiring identification of the beneficiary cannot be implemented]
\label{prop:beneficiary}
When $\iota$ is not observable, a $\pi$ conditioned on the beneficiary and the payer being
different is not $\mathcal{G}$-measurable.
\end{proposition}

\begin{proof}
The defining property of family 7 is that the party receiving $\delta$ differs from the party
paying $\pi$.
$\mathcal{P}$ identifies an identifier $j$ and $\iota(j)$ is not observable, so
$\iota(j_1)\neq\iota(j_2)$ cannot be judged on $\mathcal{G}$. By
Proposition~\ref{prop:implementable}, a $\pi$ depending on this condition cannot be written.
\end{proof}

Proposition~\ref{prop:beneficiary} does not depend on the partition: even when $\delta$ is a
state of $\mathcal{P}$, the \emph{party} receiving it cannot be identified. This is why 7-1
and 7-2 drop out in both partitions. 7-4 survives because it makes no explicit exchange and
therefore needs no identification of the beneficiary.

That mechanisms presuming parties to be distinct fail when identifiers can be created freely
is treated in mechanism design as false-name-proofness \cite{yokoo2004}.
Proposition~\ref{prop:beneficiary} is its manifestation at the level of contract form.

\section{Sieving the types}
\label{sec:protocol-filter}

The 28 types of Part~\ref{part:catalog} are passed through the sieve of the two partitions of
Table~\ref{tab:protocol-cases}.

\textbf{The criterion differs from that of Chapter~\ref{ch:solo}.} There the cut was
feasibility; here it is whether the protocol contributes. A type that receives no
contribution in the ``outside'' partition does not disappear; enforcement simply returns to a
third party. Below, a type ``receives a contribution'' when the condition governing $\pi$
enters $\mathcal{G}$ and what is enforced is a state of $\mathcal{P}$. The judgment is made
per $\Phi$, not per type (Remark~\ref{rem:case-is-phi}).

% The judgment rows exist only in the two macros below. In PDF the table does not fit the
% page height, so families 1--4 and 5--7 are split into two tables; HTML has no such
% constraint and keeps them as one. To revise a judgment, edit only the macros.
\newcommand{\protocolfilterrowsA}{%
\multicolumn{5}{@{}l}{\textbf{Family 1: one-off exchange}}\\
1-1 & immediate exchange & yes & --- & simultaneity is available only when $\delta$ and $\pi$ are the same state transition \\
1-2 & one-off advance & yes & yes & $\pi$ comes first and does not refer to $\delta$ \\
1-3 & one-off arrears & yes & --- & outside, completion of $\delta$ cannot be judged \\
1-4 & instalments & yes & partly & outside, collection of $\pi$ works but default does not reach the underlying asset \\
\midrule
\multicolumn{5}{@{}l}{\textbf{Family 2: continuing access}}\\
2-1 & flat access & yes & yes & $\pi=F$ does not refer to $\delta$ \\
2-2 & usage & yes & --- & outside, metering $\delta$ needs the party of Remark~\ref{rem:oracle} \\
2-3 & two-part tariff & yes & partly & outside, only the fixed part $F$ \\
2-4 & prepaid & yes & yes & inside, the unused balance is itself the protocol state \\
2-5 & options and warranties & yes & --- & outside, the exercise condition lies outside the protocol \\
\midrule
\multicolumn{5}{@{}l}{\textbf{Family 3: renting an asset over time}}\\
3-1 & lease and rental & yes & --- & inside only where the asset itself is a protocol state \\
3-2 & shared asset & conditionally & --- & the range in which utilization can be written as a state transition is narrow \\
3-3 & transfer plus long collection & yes & --- & inside with collateral; outside it does not reach the underlying asset \\
\midrule
\multicolumn{5}{@{}l}{\textbf{Family 4: recovery through a complement}}\\
4-1 & lock-in consumables & conditionally & --- & the family presumes the device lies outside the protocol, which fails inside \\
4-2 & freemium & yes & yes & $\kappa>0$, but it is self-issued credit and asks no enforcement of the protocol \\
4-3 & hardware subsidy plus finance & yes & --- & outside, the device lies outside the protocol \\
}
\newcommand{\protocolfilterrowsB}{%
\multicolumn{5}{@{}l}{\textbf{Family 5: outcome- and state-contingent}}\\
5-1 & success fee & conditionally & --- & only where the $y$ of $\pi=h(y)$ is a protocol quantity \\
5-2 & revenue share & yes & --- & inside, the counterparty's flow is observable on the protocol \\
5-3 & cost plus & --- & --- & actual costs lie outside the protocol either way \\
5-4 & underwriting & yes & --- & inside, $\omega$ is confined to protocol events \\
5-5 & IP licensing & yes & --- & inside, where use is observable on the protocol \\
\midrule
\multicolumn{5}{@{}l}{\textbf{Family 6: intermediation}}\\
6-1 & transaction fee & yes & --- & inside, completed in combination with the simultaneity of 1-1 \\
6-2 & escrow & yes & partly & outside, only one side is held; release needs the party of Remark~\ref{rem:oracle} \\
6-3 & spread & yes & --- & inside, inventory becomes a protocol balance \\
6-4 & charging for opportunity & conditionally & --- & only where contact can be defined as a protocol event \\
\midrule
\multicolumn{5}{@{}l}{\textbf{Family 7: separating beneficiary from payer}}\\
7-1 & advertising & --- & --- & that beneficiary and payer differ cannot be verified (Section~\ref{sec:identifier}) \\
7-2 & cross-subsidy & --- & --- & as above \\
7-3 & public payment & --- & --- & the payer sets the price and design freedom sits in the institutional layer \\
7-4 & donation and support & yes & yes & no explicit exchange, so no identification of the beneficiary is required \\
}

\ifdefined\HCode
\begin{center}
\begin{tabularx}{\textwidth}{@{}llccX@{}}
\toprule
No. & Type & Inside & Outside & Grounds \\
\midrule
\protocolfilterrowsA
\midrule
\protocolfilterrowsB
\bottomrule
\end{tabularx}
\captionof{table}{Sieving the 28 types. ``yes'' marks those that receive a contribution.}
\label{tab:protocol-filter-a}
\end{center}
\else
\begin{center}
\begin{tabularx}{\textwidth}{@{}llccX@{}}
\toprule
No. & Type & Inside & Outside & Grounds \\
\midrule
\protocolfilterrowsA
\bottomrule
\end{tabularx}
\captionof{table}{Sieving the 28 types (families 1--4). ``yes'' marks those that receive a contribution.}
\label{tab:protocol-filter-a}
\end{center}

\begin{center}
\begin{tabularx}{\textwidth}{@{}llccX@{}}
\toprule
No. & Type & Inside & Outside & Grounds \\
\midrule
\protocolfilterrowsB
\bottomrule
\end{tabularx}
\captionof{table}{Sieving the 28 types (families 5--7). ``yes'' marks those that receive a contribution.}
\label{tab:protocol-filter-b}
\end{center}
\fi

Twenty types receive a contribution in the ``inside'' partition and five in the ``outside''
partition. Where the sieve of Chapter~\ref{ch:solo} cut 28 types down to about ten, the
``inside'' sieve barely cuts at all. \textbf{This chapter's premises are a relaxation, not a
constraint.}

That 1-1 disappears in the ``outside'' partition follows from
Proposition~\ref{prop:simultaneity}: if $\delta$ lies outside the protocol, the form degrades
into either an advance or arrears.

\textbf{Receiving a contribution and being executable are, however, different things.}
The advance forms (1-2, 2-4) pass in the table for the ``inside'' partition, but executing
them requires collateral towards the counterparty and they are not open to a party without
assets (Corollary~\ref{cor:growth-needs-capital}, Corollary~\ref{cor:exclusion}).
Section~\ref{sec:collateral} treats this.

\begin{remark}[The axis along which the sieve cuts]
The five types surviving in the ``outside'' partition are 1-2, 2-1, 2-4, 4-2 and 7-4, and all
they share is that $\pi$ is not measurable with respect to $\delta$. Of the six degrees of
freedom of Section~\ref{sec:dof}, only (2) carries information for this sieve.

\textbf{This is not a proposition.} Since the ``outside'' partition was defined as the case
where $\delta$ does not enter $\mathcal{G}$, the impossibility of writing a
$\delta$-measurable $\pi$ follows from the definition. The information lies not in the claim
but in the fact that applying the sieve to 28 types individually turned out to be explained
by a single degree of freedom.
\end{remark}

\begin{remark}[The direction in which the institutional layer acts]
Types 2-4 (prepaid) and 6-2 (escrow) were the ones removed in Section~\ref{sec:bootstrap} by
deposit obligations and licensing. In the ``inside'' partition they are implemented as
protocol.

\textbf{One and the same type both drops out and returns, depending on the direction in which
the institutional layer acts.} This is an instance showing that the first form
(entry requirement) of the three in Remark~\ref{rem:institution-modes} does not act in only
one direction.
\end{remark}

\section{The collateral constraint}
\label{sec:collateral}

Section~\ref{sec:protocol-filter} showed that $\mathcal{G}$ ceases to bind. This section
identifies what does.

Consider a $\Phi$ in the ``inside'' partition with $\kappa_i>0$, that is, one in which the
operator extends credit. To enforce on non-payment of $\pi$, what is enforced must be a state
of $\mathcal{P}$. But $\mathcal{G}$ consists only of states of $\mathcal{P}$, so
\textbf{neither the counterparty's ability to pay nor their past record of performance is in
$\mathcal{G}$}.

Credit can therefore be provided for only through states locked in advance, that is,
collateral. Writing $B_i$ for the collateral one has locked for counterparty $i$ and $B_i'$
for the collateral $i$ has locked for oneself, the $\kappa_i$ of~\eqref{eq:kappa} satisfies

\begin{equation}
\max\{\kappa_i,0\} \;\leq\; B_i',
\qquad
\max\{-\kappa_i,0\} \;\leq\; B_i
\label{eq:protocol-collateral}
\end{equation}

Collateral is provided from one's own assets without duplication, so $\sum_i B_i\leq E$, and
summing~\eqref{eq:protocol-collateral} over $i$ gives the following.

\begin{proposition}[Collateral constraint]
\label{prop:collateral}
In the ``inside'' partition, the total credit that can be received does not exceed equity.
\begin{equation}
\sum_i \max\{-\kappa_i,0\} \;\leq\; E
\label{eq:protocol-finance}
\end{equation}
\end{proposition}

\begin{proof}
Summing the second inequality of~\eqref{eq:protocol-collateral} over all $i$ gives
$\sum_i\max\{-\kappa_i,0\}\leq\sum_i B_i$. The same asset cannot be locked for several
counterparties at once, so $\sum_i B_i\leq E$; combining the two gives the result.
\end{proof}

Compare~\eqref{eq:protocol-finance} with~\eqref{eq:solo-finance} of Chapter~\ref{ch:solo}.
There, under $E\approx 0$, the requirement was
$\sum\max\{-\kappa_i,0\}\geq\sum\max\{\kappa_i,0\}$: the credit drawn had to exceed the
credit extended. Here the credit that can be drawn is itself bounded by $E$.

\begin{corollary}[The substitution degenerates to one direction]
\label{cor:one-way-substitution}
In the ``inside'' partition, the substitution of Proposition~\ref{prop:substitution} becomes
one-directional: capital substitutes for credit, but credit does not substitute for capital.
\end{corollary}

\begin{proof}
By Proposition~\ref{prop:collateral} the terms with $\kappa<0$ are bounded above by $E$.
In~\eqref{eq:cash}, terms with $\kappa<0$ cannot raise funding beyond $E$. The reverse
substitution --- $E$ substituting for $\kappa<0$ --- holds as in
Proposition~\ref{prop:substitution}.
\end{proof}

Corollary~\ref{cor:one-way-substitution} is the third place in this text where
Proposition~\ref{prop:substitution} appears.

\begin{center}
\begin{tabularx}{\textwidth}{@{}lllX@{}}
\toprule
Place & $E$ & Credit & What binds there \\
\midrule
Chapter~\ref{ch:solo}, solo business & 0 & the only means of funding & the credit bootstrap problem \\
Section~\ref{sec:retiree}, those with assets & positive & built from zero & capital buys the period in which credit is built \\
This chapter, ``inside'' & the only means of funding & 0 & there is no period in which to build \\
\bottomrule
\end{tabularx}
\captionof{table}{The three places Proposition~\ref{prop:substitution} appears. The two ends line up.}
\label{tab:substitution-three-points}
\end{center}

Section~\ref{sec:credit-formation} described the route by which credit forms through
accumulated public performance. That route does not operate in the ``inside'' partition,
because $\mathcal{P}$ identifies only identifiers (Section~\ref{sec:identifier}) and cannot
attribute a past record of performance to a party.

\begin{corollary}[Exclusion of parties without assets]
\label{cor:exclusion}
When $A=\mathrm{Inv}=0$ and $E=0$, the only feasible $\Phi$ are those with $\kappa_i=0$ for
every $i$.
\end{corollary}

\begin{proof}
Substituting $E=0$ into~\eqref{eq:protocol-finance} gives $\kappa_i\geq 0$ for every $i$.
On the other hand, by Proposition~\ref{prop:substitution}, when $A=\mathrm{Inv}=E=0$ the
requirement is $\sum_i\max\{-\kappa_i,0\}\geq\sum_i\max\{\kappa_i,0\}$, whose left-hand side
is $0$, giving $\kappa_i\leq 0$.
\end{proof}

\begin{corollary}[Growth through advance payment requires capital]
\label{cor:growth-needs-capital}
In the ``inside'' partition, achieving $\CCC<0$ requires $E>0$.
\end{corollary}

\begin{proof}
By~\eqref{eq:ccc}, $\CCC=W/r$ with $W=\sum_i\kappa_i+\mathrm{Inv}$. Substituting $E=0$ into
Proposition~\ref{prop:collateral} gives $\kappa_i\geq 0$ for every $i$, which together with
$\mathrm{Inv}\geq 0$ gives $W\geq 0$ and hence $\CCC\geq 0$.
\end{proof}

Chapter~\ref{ch:solo} showed a route by which a solo business without assets creates
$\CCC<0$ through annual prepayment --- a state in which growth generates cash.
\textbf{That route closes in the ``inside'' partition}, because receiving an advance requires
collateral towards the counterparty and collateral is provided out of capital.

That is, \textbf{what remains to a party without assets is confined to types with $\kappa$
identically $0$}: 1-1 (immediate exchange), 2-2 (usage) and 6-1 (transaction fee), which
Part~\ref{part:catalog} recorded as $\kappa\approx 0$, and their composites.
In Chapter~\ref{ch:solo}, $E=0$ required $\kappa_C<0$; here it forces $\kappa=0$.
Solutions (a), (b) and (c) of Section~\ref{sec:bootstrap} all presumed a shortage of credit;
what is short here is capital.

\begin{remark}[Correspondence with existing theory]
\label{rem:collateral-priorwork}
Nothing in this section is new. The structure by which net worth constrains funding through
collateral capacity is in \cite{kiyotaki1997}, and the formulation in which collateral
constraints govern a firm's financing and risk management jointly is in
\cite{rampini2010,rampini2013}. Proposition~\ref{prop:substitution} itself has counterparts
there too (Appendix~\ref{app:priorwork}).

That the contractible space widens under premises (1) and (2) is treated in \cite{cong2019}.
The point corresponding to Corollary~\ref{cor:exclusion} --- that a mechanism demanding
collateral excludes parties without capital --- also has a literature \cite{oecd2024}.

\textbf{The role of this section is to place these on one and the same inequality,
that of~\eqref{eq:cash}.}
\end{remark}

\section{Observability}
\label{sec:protocol-observability}

Premise (2) has a secondary consequence. Because the state of $\mathcal{P}$ is verifiable by
every participant, the following quantities are individually observable in the ``inside''
partition.

\begin{center}
\begin{tabularx}{\textwidth}{@{}lXX@{}}
\toprule
Quantity & Measurement status in this text & Status in the ``inside'' partition \\
\midrule
$M(t)$ & requires per-firm data and is meaningless in aggregate & observable per identifier \\
$\kappa_j$ & as above & as above \\
$\dbar-\delta$ & only estimable from aggregates & directly observable as the unused balance \\
$B$ & this text has no corresponding quantity & observable as locked state \\
\bottomrule
\end{tabularx}
\captionof{table}{Quantities that become observable in the ``inside'' partition.}
\label{tab:protocol-observability}
\end{center}

What is observable, however, is per identifier, and by Remark~\ref{rem:additivity} the
per-party $\kappa_i$ cannot be recovered.
\textbf{The obstacle to measurement changes from the absence of aggregation to the
impossibility of it.} This differs from category (v) of Section~\ref{sec:obstacle-taxonomy}
(Remark~\ref{rem:agg-impossible}) and is a new form relative to the taxonomy of
Part~\ref{part:method}.

\subsection{A population frame becomes available}
\label{sec:protocol-frame}

It is not only quantities that become observable.

The conditioning on survival of Chapter~\ref{ch:method} arises because the objects observed
are limited to paths that did not violate~\eqref{eq:constraint}. In the ``inside'' partition
\textbf{the record of the protocol is itself the population}. Every $\Phi$ that operated
remains, including those that stopped, so no conditioning occurs.

Chapter~\ref{ch:platformfee} abandoned direct measurement of the unmanned operating period
for want of a population frame. That reason disappears here.

\subsection{What can and cannot be tested}
\label{sec:protocol-testable}

Being observable does not mean being testable. Three cases must be separated first.

\begin{center}
\begin{tabularx}{\textwidth}{@{}lX@{}}
\toprule
Case & Content \\
\midrule
Quantities the protocol enforces & observing them is not a test; it is reading the rules \\
Quantities that result from design & these are what can be tested \\
Per-party quantities & cannot be recovered (Remark~\ref{rem:additivity}) \\
\bottomrule
\end{tabularx}
\captionof{table}{Three cases among the quantities observable in the ``inside'' partition.}
\label{tab:protocol-testability}
\end{center}

The collateral constraint of Proposition~\ref{prop:collateral} belongs to the first case.
Since the protocol requires collateral, observing that collateral is posted does not test the
proposition.
\textbf{What can be tested is confined to quantities that appear as the result of an
operator's choice.}

\begin{remark}[The population differs]
\label{rem:protocol-population}
A distribution measured in the ``inside'' partition cannot be generalized to firms at large.
In the sense of Chapter~\ref{ch:method} it is a different population.
\textbf{What can be tested here is whether a deduction is correct, not whether that deduction
applies to real solo businesses.}
The propositions of Chapter~\ref{ch:solo} can be tested; nothing is said about that
chapter's objects.
\end{remark}

Chapter~\ref{ch:protocol-emp} observes these quantities in practice.

\section{Limits of this chapter}
\label{sec:protocol-limits}

\begin{enumerate}[label=(\arabic*)]
\item \textbf{This is deduction.} The conclusions depend on premises (1), (2) and (3).
The tests of Section~\ref{sec:protocol-results} measure the implications of the deduction and
do not verify the premises themselves.

\item \textbf{A rule was needed to narrow what is judged.}
Table~\ref{tab:family-assignment-order} takes $\Phi$ as given and cannot be applied directly
to ledger identifiers. Section~\ref{sec:protocol-results} places the question ``is this a
$\Phi$?'' before the assignment order (Remark~\ref{rem:is-it-phi}). That step is not part of
the rule of Chapter~\ref{ch:catalog}.

\item \textbf{Only the two ends are treated.} As Remark~\ref{rem:oracle} notes, the degree to
which the third party is removed in the ``outside'' partition is continuous. This chapter
treats only full removal and no removal, not the middle. Most existing arrangements are in
the middle.

\item \textbf{A limit of the definitions remains.} The unobservability of $\iota$ in
Section~\ref{sec:identifier} originates in fixing $N$ in Chapter~\ref{sec:phi}. This chapter
records it and does not extend the domain. It is the same place as the circulation of claims
in Section~\ref{sec:token-limit}.

\item \textbf{$\phi$ is not treated.} This chapter argues only about the feasibility of
$\Phi$ and does not apply the three-way decomposition of surplus of
Chapter~\ref{sec:surplus}. An argument of the same form as the remark in
Chapter~\ref{part:unmanned} --- that the replicability of a protocol puts pressure on
$\phi_{\rm prod}$ --- could hold, but judging it requires measured fees and is not attempted
here.
\end{enumerate}

%============================================================
